\section{Optimization problem}\label{sec:sp:opt}

\subsection{The planner's problem}\label{sec:sp:opt:problem}

\smalltitle{Controls and states.} The planner chooses paths for the eight controls
\begin{align}
    \left\{C,\;I,\;N,\;D,\;\varpi,\;K^R,\;W,\;\Omega\right\},
\label{eq:sp:controls}
\end{align}
which determine the remaining flows through the definitions
\begin{align}
R^R &= \mathcal R\left(\varpi W,K^R\right),
&
R &= N+R^R,
&
Y &= F\left(K-K^R,\,R,\,P\right),
&
G &= I+\delta K,
\label{eq:sp:definitions}
\end{align}
and it carries the five states
\begin{align}
\left\{K,\;S,\;X,\;P,\;M^K\right\},
\label{eq:sp:states}
\end{align}
with $K_0,S_0,X_0,P_0,M^K_0$ given.

Two of the eight controls are not free. Given the remaining six, the ledger~\eqref{eq:sp:ledger} and the intensity condition~\eqref{eq:sp:intensity} are two equations in $W$ and $\Omega$, and the goods constraint~\eqref{eq:sp:goods} removes one more degree of freedom, leaving five independent instruments against the five states. We nonetheless carry $W$ and $\Omega$ as controls with the two relations adjoined, so that their multipliers $p^W$ and $\zeta$ appear explicitly and can be read as shadow prices.


\begin{problem}[The planner's problem]\label{prob:sp}
Choose $\left\{C,I,N,D,\varpi,K^R,W,\Omega\right\}_{t\ge0}$ to maximise

\noindent $U_0=\int_0^\infty\left[u(C)-v(P)\right]e^{-\rho t}dt$ subject to, at every $t$,
\begin{align}
&\text{\emph{goods:}}
&
F\left(K-K^R,R,P\right) &= C+I+\delta K+C^N(N,S)+C^D(D,X)+c^W\!\left(\varpi\right)W,
\label{eq:sp:goods}
\\[0.2em]
&\text{\emph{ledger:}}
&
W &= R-\Omega\phi^I G+\delta M^K+\mathcal N,
\label{eq:sp:ledger}
\\[0.2em]
&\text{\emph{intensity:}}
&
R &= \Omega\left(\mathcal D+\phi^IG\right),
\label{eq:sp:intensity}
\end{align}
and to the laws of motion
\begin{align}
\dot K &= I,
&
\dot S &= D-N,
&
\dot X &= D,
&
\dot P &= \Xi-\theta(P)P,
&
\dot M^K &= \Omega\phi^I G-\delta M^K,
\label{eq:sp:states-motion}
\end{align}
with waste bases $\mathcal{N}$ and $\mathcal{D}$ defined in \eqref{eq:sp:setup:base}, $c^W \equiv c^c+ c^T(\varpi)$, the definitions in~\eqref{eq:sp:definitions}, the emission flow from \eqref{eq:sp:setup:Xi} and domains $C,I+\delta K,N,D,K^R,W,\Omega\ge0$, $\varpi\in\left[0,1\right]$.
\end{problem}

The two static relations adjoined alongside the goods constraint are the accounting identities of Section~\ref{sec:sp:setup:aggregatematerial}, restated in the variables the planner chooses. Equation~\eqref{eq:sp:ledger} is the ledger~\eqref{eq:sp:setup:ledger}: waste generated equals the material entering production, less the net accumulation of material embodied in the capital stock, plus the mass displaced from nature. Equation~\eqref{eq:sp:intensity} is the balance of embodied materials in the final good, \eqref{eq:sp:setup:ybalance_OmegaPhi}, and is what determines $\Omega$. The two are distinct restrictions rather than two forms of one identity: substituting the second into the first returns
\begin{align}
W &= \Omega\mathcal D+\delta M^K+\mathcal N,
\label{eq:sp:ledger-streams}
\end{align}
which is the sum of the individual waste streams~\eqref{eq:sp:setup:totalwaste_generic}. The pair thus carries no content beyond the process-level accounting, but together they pin down $W$ and $\Omega$, as noted above.

$\Omega$ itself has a somewhat complicated interpretation, as it captures the level of the average material intensity of production while also adjusting to reflect whether final goods are put to more or less material-intensive purposes. From~\eqref{eq:sp:intensity}, $\Omega=R/\left(\mathcal D+\phi^IG\right)$ scales with $R$ for a fixed structure of uses, but for a given $R$ it falls whenever $Y$ is shifted towards a use carrying a larger coefficient ($\omega^j$ or $\phi^I$). We refer to the relative sizes of the different uses of the final good as the \textit{composition of activities/uses} of the good.

The channel through which that composition reaches the allocation is narrow, and worth marking before the optimality conditions. Given $R$, the ledger says that everything entering production becomes waste except what gross investment stores in $M^K$. The composition therefore matters only by moving $\Omega$ through $\Phi^I=\Omega\phi^I$ i.e. the quantity of material held back from today's waste flow. This is precisely what the multiplier $\zeta$ on~\eqref{eq:sp:intensity} prices, and Section~\ref{sec:sp:opt:foc} shows that the channel closes entirely when $\phi^I=0$.


Next, we note that $\Xi$ is written as linear in $W$ and $R^R$ separately, so that no derivative of $a(\cdot)$ ever appears inside a pollution term. All the curvature of the recycling
technology enters through $\mathcal R$ and through $\alpha$ and $a'$ alone.\footnote{Recall
from~\eqref{eq:sp:setup:recycling-derivatives} that $\mathcal R_T=\alpha\equiv a-a'x$ is the marginal
recycling yield and $\mathcal R_{K^R}=a'$.}



\smalltitle{The Lagrangian.} Write the current-value Hamiltonian with the static constraints
adjoined. It is convenient to normalise every multiplier by $\lambda$, the multiplier on the goods
constraint~\eqref{eq:sp:goods}, so that all shadow prices below are measured in units of the final
good rather than in utils. Accordingly, let the current-value costates be
\begin{align}
\underbrace{\lambda q}_{K},
&&
\underbrace{\lambda p^S}_{S},
&&
\underbrace{\lambda p^X}_{X},
&&
\underbrace{-\lambda p^P}_{P},
&&
\underbrace{-\lambda p^M}_{M^K},
\label{eq:sp:costates-def}
\end{align}
the two minus signs anticipating that both $P$ and $M^K$ are likely liabilities, so that $p^P$ and $p^M$ read as shadow \emph{costs}. Let $\lambda p^W$ be the multiplier on the ledger~\eqref{eq:sp:ledger} written with $W$ on the left, and $\lambda\zeta$ the multiplier on the intensity definition~\eqref{eq:sp:intensity} written with
$R$ on the left. Then
\begin{align}
\mathcal L
={}& u(C)-v(P)
\notag\\
&+\lambda\Big[F\left(K-K^R,R,P\right)-C-I-\delta K-C^N(N,S)-C^D(D,X)-c^W\!\left(\varpi\right)W\Big]
\notag\\
&+\lambda q\,I
\;+\;\lambda p^S\left(D-N\right)
\;+\;\lambda p^X D
\;-\;\lambda p^P\left[\Xi-\theta(P)P\right]
\;-\;\lambda p^M\left[\Omega\phi^IG-\delta M^K\right]
\notag\\
&+\lambda p^W\Big[W-N-R^R+\Omega\phi^IG-\delta M^K-\mathcal N\Big]
\;+\;\lambda \zeta\Big[N+R^R-\Omega\left(\mathcal D+\phi^IG\right)\Big],
\label{eq:sp:lagrangian}
\end{align}
with relevant definitions substituted throughout.\footnote{This includes $R=N+R^R$, $R^R=\mathcal R\left(\varpi W,K^R\right)$, $G=I+\delta K$, $Y=F\left(K-K^R,R,P\right)$
and $C^W=c^W\!\left(\varpi\right)W$ substituted throughout.} Here, the multiplier $p^W$ measures the \textit{social cost} of waste in units of final goods and $\zeta$ measures the social value of embodied materials in units of final goods.\footnote{To see this, assume that we could increase waste from nature $\mathcal{N}$ marginally without changing any other controls, i.e. if more waste was generated out of the blue; we would have $\partial \mathcal L^*/\partial \mathcal N = - \lambda p^W$ which would be negative (i.e. a cost) exactly if $p^W>0$. Similarly, perturb the condition $R + \kappa = \Omega \left(\mathcal{D}+\phi^IG\right)$ such that we could increase embodied materials (right hand side) marginally without increasing the amount of materials $R$ that goes into the production process; we would then have $\partial \mathcal L^*/\partial \kappa = \lambda \zeta$ which would be positive if $\zeta>0$.} 



\subsection{Optimality conditions}\label{sec:sp:opt:foc}

The conditions are collected here; Appendix~\ref{sec:app:sp} derives each of them, decomposes it into the separate margins it balances, and checks the units. Define
\begin{align}
\Psi &\equiv F_R\left(1-\zeta\Omega^Y\right)+\zeta,
&
\sigma &\equiv \frac{\Phi^IG}{R}\;\in\left[0,1\right],
&
r &\equiv \rho-\frac{\dot\lambda}{\lambda},
\label{eq:sp:shorthand}
\end{align}
where $\Psi$ is the value of a tonne of raw material delivered to production and $\sigma$ is the share of material embodied in new capital. $\Psi$ captures output net of embodied material output requires plus relaxation of the intensity condition, but excludes the ledger: the tonne must eventually leave as waste, and that is priced separately by $p^W$.

\smalltitle{Controls.} Each condition is stated below with the margin it balances. Appendix~\ref{sec:app:sp} derives each one, decomposes it term by term, and checks the units.

\emph{Material intensity.} Raising $\Omega$ stores material in the capital stock rather than releasing it now. The gain is the price gap $p^W-p^M$, scaled by the share $\sigma$ of material that storage actually captures, and the shadow price of embodied material must equal it:
\begin{align}
\zeta &= \sigma\left(p^W-p^M\right).
\label{eq:sp:sum:Omega}
\end{align}

\emph{Consumption.} A unit consumed absorbs one final good and carries $\Omega^C$ tonnes of embodied material away from the storage margin. Marginal utility therefore exceeds $\lambda$ by that wedge, so $\lambda$ is not the marginal utility of consumption in this model:
\begin{align}
u'\left(C\right)/\lambda &= 1+\zeta\Omega^C .
\label{eq:sp:sum:C}
\end{align}

\emph{Investment.} A unit invested buys a unit of capital and stores $\Phi^I$ tonnes, of which the fraction $\sigma$ is clawed back because embodying material also consumes material drawn through $R$. The shadow price of installed capital departs from unity by exactly the net storage value:
\begin{align}
q &= 1-\Phi^I\left(1-\sigma\right)\left(p^W-p^M\right).
\label{eq:sp:sum:I}
\end{align}

\emph{Extraction.} A tonne extracted is worth $\Psi$ on delivery. It must cover the extraction effort, grossed up for the embodied material the effort itself carries, the in-situ rent, and the waste it becomes (both the tonne and the overburden displaced alongside it):
\begin{align}
\Psi &= C^N_N\left(1+\zeta\Omega^N\right)+p^S+p^W\left(1+\Omega^{N,S}\right).
\label{eq:sp:sum:N}
\end{align}

\emph{Exploration.} A tonne discovered creates the rent $p^S$ and uses up a deposit, $p^X<0$. Marginal costs cover exploration grossed up by the embodied material itself carries and from waste generation:  
\begin{align}
p^S+p^X &= C^D_D\left(1+\zeta\Omega^D\right)+p^W\Omega^{D,S} .
\label{eq:sp:sum:D}
\end{align}

\emph{Waste.} A tonne of waste absorbs collection and treatment, adds to the emission flow, and returns $\alpha\varpi$ tonnes; the factor $1-\alpha\varpi>0$ on the left is the net tonnage the ledger attributes to it once the recycled part is returned to $R$:
\begin{align}
p^W\left(1-\alpha\varpi\right) &= c^W\left(1+\zeta\Omega^W\right)+p^P\left[1-\varpi\left(1-d^W\right)-d^W\alpha\varpi\right]-\alpha\varpi\,\Psi .
\label{eq:sp:sum:W}
\end{align}

\smalltitle{The two bounded controls.} The six conditions above are equalities because the controls
they price are interior at any candidate optimum: consumption and extraction by the Inada conditions,
investment and the two accounting controls because their bounds are never approached. The remaining
two controls are the recycling block, and they are different in kind. Treatment is bounded on both
sides, $\varpi\in\left[0,1\right]$, and recycling capital is bounded below, $K^R\ge0$; both bounds are
reachable, and Section~\ref{sec:sp:opt:phases} shows that one of them is in fact binding over an
initial interval of the optimal path. Their conditions are therefore stated as Kuhn--Tucker systems:
each corner is reported together with the inequality that makes it optimal, and the interior equality
together with the inequality that makes \emph{it} optimal, so that the three branches partition the
possibilities. To keep the branches readable, write the marginal gain on each of the two margins as
\begin{align}
\mathcal G^{\varpi} &\equiv \alpha\left(\Psi-p^W\right)+p^P\left[\left(1-d^W\right)+d^W\alpha\right],
&
\mathcal G^{K} &\equiv \Psi-p^W+d^Wp^P,
\label{eq:sp:sum:gains}
\end{align}
both in final goods per tonne: $\mathcal G^\varpi$ is the value of routing one more tonne of waste to
treatment, and $\mathcal G^{K}$ the value of one more tonne of secondary material recovered by
capital, so that the gain from a unit of recycling capital is $a'\!\left(x\right)\mathcal G^K$.

\emph{Treatment share.} Treating one more tonne yields $\alpha$ tonnes of secondary material, each
worth $\Psi-p^W$ rather than $\Psi$ because the recovered material must itself leave as waste later,
and it diverts the tonne from the environment. The gain $\mathcal G^\varpi$ is set against a rising
marginal treatment cost, and since $c^{T\prime}\left(0\right)=0$ by~\eqref{eq:sp:setup:waste-cost} the
lower corner requires the gain itself to be non-positive:
\begin{subequations}
\label{eq:sp:sum:varpi}
\begin{align}
\varpi &= 0,
&&\text{if}\quad \mathcal G^{\varpi} \;\le\; 0,
\label{eq:sp:sum:varpi-lower}
\\
\mathcal G^{\varpi} &= c^{T\prime}\left(\varpi\right)\left(1+\zeta\Omega^W\right),
&&\text{if}\quad 0 \;<\; \mathcal G^{\varpi} \;<\; c^{T\prime}\left(1\right)\left(1+\zeta\Omega^W\right),
\label{eq:sp:sum:varpi-interior}
\\
\varpi &= 1,
&&\text{if}\quad \mathcal G^{\varpi} \;\ge\; c^{T\prime}\left(1\right)\left(1+\zeta\Omega^W\right).
\label{eq:sp:sum:varpi-upper}
\end{align}
\end{subequations}
Convexity of $c^T$ makes the marginal cost strictly increasing, so the three branches are exhaustive
and mutually exclusive, and the interior root is unique when it exists. Part~(i) of
Proposition~\ref{prop:sp:phases} shows that~\eqref{eq:sp:sum:varpi-lower} is never satisfied: the
pollution term $p^P\left(1-d^W\right)$ in $\mathcal G^\varpi$ is strictly positive whatever the
recycling margin is doing, so the marginal gain from the first unit of treatment is strictly positive
while its marginal cost is zero. The upper corner is a genuine possibility and is not ruled out by any
assumption made so far; whether it is ever reached depends on how steeply $c^T$ rises as
$\varpi\to1$.\footnote{If one wants full treatment to be excluded by technology rather than by
parameter values, it is enough to impose $\lim_{\varpi\to1}c^{T\prime}\left(\varpi\right)=\infty$.}

\emph{Recycling capital.} A final good moved into recycling yields $a'$ tonnes, valued at $\mathcal G^K$. It costs the marginal product forgone in production, net
of the material that output would have required. The Kuhn-Tucker conditions then read:
\begin{subequations}
\label{eq:sp:sum:KR}
\begin{align}
K^R &= 0,
&&\text{if}\quad a'\!\left(0\right)\mathcal G^{K} \;\le\; F_K\left(1-\zeta\Omega^Y\right),
\label{eq:sp:sum:KR-corner}
\\
a'\!\left(x\right)\mathcal G^{K} &= F_K\left(1-\zeta\Omega^Y\right),
&&\text{if}\quad a'\!\left(0\right)\mathcal G^{K} \;>\; F_K\left(1-\zeta\Omega^Y\right).
\label{eq:sp:sum:KR-interior}
\end{align}
\end{subequations}
The upper bound $K^R\le K$ is
never stated because it never binds: $K^Y\to0$ sends $F_K\to\infty$ and
violates~\eqref{eq:sp:sum:KR-interior} long before it is reached.

\emph{The two recycling conditions restated.} Whenever extraction is interior, the combination
$\Psi-p^W$ inside $\mathcal G^\varpi$ and $\mathcal G^K$ can be eliminated using the extraction
condition~\eqref{eq:sp:sum:N}, and the unit waste charge cancels (see appendix \ref{sec:app:sp}). Define the \emph{virgin displacement value}
\begin{align}
V &\equiv \Psi-p^W+d^Wp^P \;=\; C^N_N\left(1+\zeta\Omega^N\right)+p^S+\Omega^{N,S}p^W+d^Wp^P,
\label{eq:sp:sum:V}
\end{align}
we have $\mathcal G^K=V$ and $\mathcal G^\varpi=\alpha V+p^P\left(1-d^W\right)$, so that the two
Kuhn--Tucker systems read
\begin{subequations}
\label{eq:sp:sum:recycling-virgin}
\begin{align}
\varpi &= 0
&&\text{if}\quad \alpha V+p^P\left(1-d^W\right)\le 0,
\notag
\\
\alpha V+p^P\left(1-d^W\right) &= c^{T\prime}\left(\varpi\right)\left(1+\zeta\Omega^W\right)
&&\text{if}\quad 0<\alpha V+p^P\left(1-d^W\right)<c^{T\prime}\left(1\right)\left(1+\zeta\Omega^W\right),
\notag
\\
\varpi &= 1
&&\text{if}\quad \alpha V+p^P\left(1-d^W\right)\ge c^{T\prime}\left(1\right)\left(1+\zeta\Omega^W\right),
\label{eq:sp:sum:recycling-virgin-varpi}
\end{align}
and
\begin{align}
K^R &= 0
&&\text{if}\quad a'\!\left(0\right)V\le F_K\left(1-\zeta\Omega^Y\right),
\notag
\\
a'\!\left(x\right)V &= F_K\left(1-\zeta\Omega^Y\right)
&&\text{if}\quad a'\!\left(0\right)V> F_K\left(1-\zeta\Omega^Y\right).
\label{eq:sp:sum:recycling-virgin-KR}
\end{align}
\end{subequations}
\emph{Recycling is paid the cost of the virgin tonne it displaces, and nothing else}: extraction
effort with the material it carries, the in-situ rent, the overburden that tonne would have brought
up, and the leakage it would have suffered. The charge $p^W$ is absent because the recovered tonne and
the virgin tonne it replaces face the same charge on the way out; recycling changes the route a
tonne takes into production, never the fact that it must eventually leave. The one motive $V$ cannot
price is the diversion of the untreated tonne from the environment, $p^P\left(1-d^W\right)$, which is
why that term survives separately on the treatment margin and has no counterpart on the capital
margin.

\smalltitle{Costates.} Each state carries the costate named in~\eqref{eq:sp:costates-def}, with the transversality condition $\lim_{t\to\infty}e^{-\rho t}\mu Z=0$. That condition is what selects the forward solution reported alongside each law of motion below, and it gives every price the same asset-pricing form: the discounted stream of the dividend the stock pays while it is held.

\emph{Capital.} A unit of capital pays its marginal product, net of the embodied material the extra output requires, and runs off at rate $\delta$:
\begin{align}
\dot q &= \left(r+\delta\right)q-F_K\left(1-\zeta\Omega^Y\right),
\qquad
q\left(t\right)=\int_t^\infty F_K\left(1-\zeta\Omega^Y\right)e^{-\int_t^s\left(r+\delta\right)d\tau}ds .
\label{eq:sp:sum:K}
\end{align}
Dividing the first expression by $q$ returns the familiar arbitrage form $r=F_K\left(1-\zeta\Omega^Y\right)/q-\delta+\dot q/q$.\footnote{The storage motive of~\eqref{eq:sp:sum:I} appears nowhere in it: depreciation destroys capital and releases its embodied material at the same rate, so the two cancel in the return.}

\emph{Reserves.} A tonne left in the ground pays a dividend by making the extraction still to come cheaper, since $C^N_S<0$:
\begin{align}
\dot p^S &= rp^S-\left|C^N_S\right|\left(1+\zeta\Omega^N\right),
\qquad
p^S\left(t\right)=\int_t^\infty\left|C^N_S\right|\left(1+\zeta\Omega^N\right)e^{-\int_t^s r\,d\tau}ds .
\label{eq:sp:sum:S}
\end{align}
This is Hotelling's rule with a stock effect: the rent is positive and need not grow at the rate of interest, because the reserve yields a flow benefit as well as a capital gain.

\emph{Discoveries.} Cumulative discovery is a pure liability, raising the cost of all future exploration:
\begin{align}
\dot p^X &= rp^X+C^D_X\left(1+\zeta\Omega^D\right),
\qquad
p^X\left(t\right)=-\int_t^\infty C^D_X\left(1+\zeta\Omega^D\right)e^{-\int_t^s r\,d\tau}ds .
\label{eq:sp:sum:X}
\end{align}
Its negative sign is what stops exploration from being a costless substitute for extraction in~\eqref{eq:sp:sum:D}.

\emph{Pollution.} The stock damages utility and output, the two together giving a marginal damage $\mathcal V\equiv v'\left(P\right)/\lambda-F_P\left(1-\zeta\Omega^Y\right)>0$ in final goods per tonne, and it decays at the \emph{marginal} rate $\theta+\theta'P$ rather than the average rate $\theta$:
\begin{align}
\dot p^P &= \left[r+\theta+\theta'P\right]p^P-\mathcal V,
\qquad
p^P\left(t\right)=\int_t^\infty\mathcal V\,e^{-\int_t^s\left[r+\theta+\theta'P\right]d\tau}ds .
\label{eq:sp:sum:P}
\end{align}
Both damages are costs, so $p^P>0$ unconditionally. Since $\theta'<0$ the marginal decay rate is below the average, and if it falls far enough that $r+\theta+\theta'P\le0$ the integral need not converge (this model's version of a tipping point).

\emph{Embodied materials.} A tonne held in the capital stock is released at rate $\delta$, becoming waste when it is:
\begin{align}
\dot p^M &= \left(r+\delta\right)p^M-\delta p^W,
\qquad
p^M\left(t\right)=\int_t^\infty\delta\,p^W\!\left(s\right)e^{-\int_t^s\left(r+\delta\right)d\tau}ds .
\label{eq:sp:sum:M}
\end{align}
The integrated form is the one to keep in view: $p^M$ is a \emph{forward} average of $p^W$, not a function of its current value.

\emph{Consumption.} Combining~\eqref{eq:sp:sum:C} with the definition of $r$ gives the Keynes--Ramsey rule with one extra term, where $\eta\equiv-u''C/u'$:
\begin{align}
\eta\,\frac{\dot C}{C} &= r-\rho-\frac{d}{dt}\ln\left(1+\zeta\Omega^C\right).
\label{eq:sp:sum:Cgrowth}
\end{align}
A consumption wedge that is constant over time affects only the level of consumption; it is a wedge that is \emph{changing} that tilts the path.



\smalltitle{Reading the system.} Three prices carry the materials side of the problem, and every condition above is built from them together with the standard terms. $\Psi$ is what a tonne is worth entering production; $p^W$ is what it costs on the way out; and $\zeta$ is what it is worth to have a tonne stored in the capital stock rather than diffused. The separation is deliberate. A tonne is valued once when it arrives, charged once when it leaves, and the difference between the two dates is priced by $\zeta$.

Relative to the same model without materials accounting, two patterns recur. First, every flow measured in final goods is grossed up by a dimensionless factor $1\pm\zeta\Omega^j$, because the activity carries embodied material along with the resources it absorbs: consumption in~\eqref{eq:sp:sum:C}, extraction and exploration effort in~\eqref{eq:sp:sum:N}--\eqref{eq:sp:sum:D}, treatment in~\eqref{eq:sp:sum:varpi}, and output itself wherever $F_K$ or $F_R$ appears. These wedges are what the composition of activities/uses does to the allocation, and they vanish with $\zeta$. Second, the ledger charges each tonne entering $R$ exactly once, whatever route it took: this is why extraction pays $p^W$ per tonne in~\eqref{eq:sp:sum:N} while exploration pays only for the mass it displaces in~\eqref{eq:sp:sum:D}, and why the recovered tonne in~\eqref{eq:sp:sum:varpi} is worth $\Psi-p^W$ rather than $\Psi$.

The system closes in the usual way. Given the five states and the five costates $\left(q,p^S,p^X,p^P,p^M\right)$, the eight first-order conditions~\eqref{eq:sp:sum:Omega}--\eqref{eq:sp:sum:KR} together with the three static constraints of Problem~\ref{prob:sp} form eleven equations in eleven unknowns: the eight controls and the three prices $\lambda$, $p^W$ and $\zeta$ that have no state of their own. Two of the eleven are equations only branch by branch: on the treatment and recycling-capital margins the system selects between a corner and an interior root according to~\eqref{eq:sp:sum:varpi} and~\eqref{eq:sp:sum:KR}, which is one equation either way and is what makes the count work in both phases. The costate equations~\eqref{eq:sp:sum:K}--\eqref{eq:sp:sum:M} and the laws of motion~\eqref{eq:sp:states-motion} then propagate the ten remaining objects, with the consumption rule~\eqref{eq:sp:sum:Cgrowth} following from~\eqref{eq:sp:sum:C} and the definition of $r$ rather than adding a restriction of its own.

\smalltitle{Corners outside the recycling block.} Two further non-negativity constraints deserve a
word, because although neither is reachable on the margins the model is built to study, one of them
does bind eventually. Extraction is interior at every date on which materials are essential: $N=0$
with $R^R=0$ would give $F_R=\infty$, and the argument in the proof of
Proposition~\ref{prop:sp:phases} turns this into $N>0$ wherever recycling is shut down. Once recycling
is running, however, $R^R>0$ keeps $F_R$ finite and $N=0$ becomes admissible, with
\begin{align}
N &= 0
\quad\text{if}\quad
\Psi \;\le\; C^N_N\left(0,S\right)\left(1+\zeta\Omega^N\right)+p^S+p^W\left(1+\Omega^{N,S}\right),
\label{eq:sp:sum:N-corner}
\end{align}
and~\eqref{eq:sp:sum:N} holding with equality otherwise. Exploration is the one corner that is
certain to bind: $p^X$ accumulates the cost of past discoveries without bound
by~\eqref{eq:sp:sum:X}, so the net value $p^S+p^X$ of a discovery must eventually fall below its
marginal cost and
\begin{align}
D &= 0
\quad\text{if}\quad
p^S+p^X \;\le\; C^D_D\left(0,X\right)\left(1+\zeta\Omega^D\right)+p^W\Omega^{D,S},
\label{eq:sp:sum:D-corner}
\end{align}
with~\eqref{eq:sp:sum:D} holding with equality otherwise. Section~\ref{sec:sp:opt:results} uses
exactly this: exploration ceases at a finite date, after which $\dot S=-N<0$ and the reserve is run
down. We state these two in passing rather than in the branch-by-branch form used for the recycling
block, because neither carries any of the paper's argument --- the phase structure of
Section~\ref{sec:sp:opt:phases} is about $\varpi$ and $K^R$ alone --- but both have to be carried in
the numerical implementation.

\subsection{Phases of the optimal path}\label{sec:sp:opt:phases}

The Kuhn--Tucker systems~\eqref{eq:sp:sum:varpi} and~\eqref{eq:sp:sum:KR} left three corners open:
$\varpi=0$, $\varpi=1$ and $K^R=0$. This section settles which of them can be optimal. They do not
behave alike. The lower treatment corner can never be optimal; the recycling corner can, and when it
is, the economy is running a materially \emph{linear} allocation by choice rather than by assumption.
This matters for the paper's argument as much as for the numerical implementation: the model does not
build circularity in, it produces it, and it dates the transition endogenously. The upper treatment
corner $\varpi=1$ is left where~\eqref{eq:sp:sum:varpi-upper} leaves it, as a possibility governed by
the calibration of $c^T$; nothing below depends on excluding it.

It is worth being clear about which corners are even candidates. Because the recycling technology is
constant returns with both inputs essential --- $\mathcal R\left(T,0\right)=a(0)T=0$ and
$\mathcal R\left(0,K^R\right)=\lim_{x\to\infty}a\left(x\right)K^R/x=0$ --- setting either argument to
zero shuts recycling down entirely, so $R^R=0$ in both cases. The configuration $\varpi=0$ with
$K^R>0$ is therefore never worth considering separately: capital with no waste to work on is idle, and
part~(i) below rules the treatment corner out in any case.

\begin{proposition}[Phases of the optimal path]\label{prop:sp:phases}
Suppose materials are essential in production, $\lim_{R\to0}F_R=\infty$, and that the material wedges
are regular in the sense that $1+\zeta\Omega^j>0$ for $j\in\left\{N,W\right\}$ and
$1-\zeta\Omega^Y>0$. Then along any optimal path:
\begin{enumerate}[label=\emph{(\roman*)},leftmargin=2.6em]
\item \emph{Treatment is always positive.} $\varpi>0$ at every date, so
branch~\eqref{eq:sp:sum:varpi-lower} is empty. The pure ``linear'' allocation, in which waste is
neither treated nor recycled, is never optimal.

\item \emph{At a recycling corner, extraction is positive and waste is a liability.} At any date at
which $K^R=0$ we have $N>0$, so the extraction condition~\eqref{eq:sp:sum:N} holds with equality, and
\begin{align}
p^W &= c^W\left(1+\zeta\Omega^W\right)+p^P\left[1-\varpi\left(1-d^W\right)\right] \;>\;0,
&
V &\ge C^N_N+p^S+d^Wp^P\;>\;0 .
\label{eq:sp:phases:corner-prices}
\end{align}
In particular the sign ambiguity in $p^W$ is resolved in this phase: waste cannot be a resource where
nothing is recovered from it.

\item \emph{The recycling corner binds on a threshold.} $K^R=0$ is optimal at $t$ if and only if
\begin{align}
a'\!\left(0\right) V &\le F_K\left(1-\zeta\Omega^Y\right),
\label{eq:sp:phases:threshold}
\end{align}
and by~\eqref{eq:sp:setup:recycling} and part~(ii) both sides are finite and strictly positive, so the
condition is neither vacuous nor a knife-edge. The economy therefore runs through at most two phases:
\begin{center}
\begin{tabular}{llll}
\toprule
& $\varpi$ & $K^R$ & materials \\
\midrule
\emph{Semi-linear} & $>0$ & $=0$ & $a=0$, $R=N$, circularity multiplier $=1$, $p^W>0$ \\
\emph{Circular} & $>0$ & $>0$ & multiplier $>1$, $p^W$ of either sign \\
\bottomrule
\end{tabular}
\end{center}
In the semi-linear phase treatment is chosen for the pollution motive alone, $p^P\left(1-d^W\right)=
c^{T\prime}\left(\varpi\right)\left(1+\zeta\Omega^W\right)$, and $\varpi$ is the share of waste placed
under controlled disposal rather than the share recycled.

\item \emph{The switch is a kink, not a jump.} At any date $t_R$ at which~\eqref{eq:sp:phases:threshold}
holds with equality, the paths of $x$, $K^R$, $\alpha$ and $\varpi$ are continuous, with
$x\left(t_R\right)=K^R\left(t_R\right)=\alpha\left(t_R\right)=0$; the discontinuity is in $\dot K^R$
and $\dot\varpi$ only.
\end{enumerate}
\end{proposition}

\begin{proof}
\emph{A preliminary step covering both corners.} Suppose $R^R=0$ at some date, which by the
essentiality of both recycling inputs is what either corner implies. Then $R=N$, and $N=0$ is not
optimal: it would give $F_R=\infty$ and hence $\Psi=\infty$ in~\eqref{eq:sp:sum:N} against a finite
right-hand side. So $N>0$ and~\eqref{eq:sp:sum:N} holds with equality, which is what licenses the
identity~\eqref{eq:sp:sum:V} for $V$. Next, $R^R=0$ implies $\alpha\varpi=0$: at $K^R=0$ we have
$\alpha=0$ by the argument following~\eqref{eq:sp:setup:recycling-derivatives}, and at $\varpi=0$ the
product vanishes directly. The waste condition~\eqref{eq:sp:sum:W} then collapses to the expression
in~\eqref{eq:sp:phases:corner-prices}, whose bracket is bounded below by $d^W>0$, so $p^W>0$. Since
$p^S>0$ by~\eqref{eq:sp:sum:S} and $p^P>0$ by~\eqref{eq:sp:sum:P}, and $\Omega^{N,S}\ge0$, every term
of $V$ in~\eqref{eq:sp:sum:V} is then positive, giving the stated bound.

\emph{(i)} The lower corner branch~\eqref{eq:sp:sum:varpi-lower} requires
$\mathcal G^\varpi\le0$, which in the form~\eqref{eq:sp:sum:recycling-virgin-varpi} reads
$\alpha V+p^P\left(1-d^W\right)\le c^{T\prime}\left(0\right)\left(1+\zeta\Omega^W\right)$. The
right-hand side is zero by~\eqref{eq:sp:setup:waste-cost}. On the left, $\alpha\in\left[0,1\right]$ and
$V>0$ by the preliminary step, so $\alpha V\ge0$, while $p^P\left(1-d^W\right)>0$ by $d^W<1$
and~\eqref{eq:sp:sum:P}. The inequality fails, so $\varpi>0$. Note that the argument does not require
$\alpha$ to be pinned down, which matters because $x=K^R/\varpi W$ is undefined when $\varpi$ and
$K^R$ vanish together.

\emph{(ii)} Immediate from the preliminary step, using $\varpi>0$ from~(i) so that $x=0$ and
$\alpha=0$ are well defined at the corner.

\emph{(iii)} This is the branching inequality
of~\eqref{eq:sp:sum:recycling-virgin-KR}, restated here as an equivalence rather than as a
sufficient condition. Necessity is the corner branch itself; sufficiency follows because $a''<0$
makes $a'\left(x\right)V$ strictly decreasing in $x$, so if the corner inequality fails the interior
root exists and is unique, and if it holds no interior root does. Finiteness of the left-hand side is
$a'\left(0^+\right)<\infty$ from~\eqref{eq:sp:setup:recycling}, strict positivity is $V>0$ from~(ii);
the right-hand side is positive by the regularity assumption. The characterisation of the semi-linear
phase follows by setting $\alpha=0$ in~\eqref{eq:sp:sum:recycling-virgin-varpi} and $a=0$
in~\eqref{eq:sp:setup:ledger-collected}.

\emph{(iv)} Since $a'$ is continuous and strictly decreasing with $a'\left(0^+\right)$ finite, the
interior condition $a'\left(x\right)V=F_K\left(1-\zeta\Omega^Y\right)$ defines $x$ as a continuous
function of the ratio $F_K\left(1-\zeta\Omega^Y\right)/V$ on the region where that ratio is below
$a'\left(0^+\right)$, and $x\to0$ as the ratio approaches $a'\left(0^+\right)$ from below. Hence $x$,
and with it $K^R=x\varpi W$ and $\alpha=a-a'x$, is continuous at $t_R$ with value zero. Because
$\alpha$ is continuous and vanishes at $t_R$, the treatment
condition~\eqref{eq:sp:sum:recycling-virgin-varpi} is continuous there too, so $\varpi$ is continuous. The
slopes differ on the two sides because $\alpha V$ is identically zero before $t_R$ and rises after it.
\end{proof}

\smalltitle{What moves the threshold.} Condition~\eqref{eq:sp:phases:threshold} compares the return on
a final good placed in recycling, $a'\left(0^+\right)V$, with its return in production,
$F_K\left(1-\zeta\Omega^Y\right)$. It therefore binds where capital is scarce and virgin material is
cheap: a low $C^N_N$ and a low rent $p^S$, which is to say abundant and accessible reserves, together
with a high marginal product of capital. That is a description of an early stage of development, and
the phase is self-terminating under depletion alone: as $S$ falls both $C^N_N$ and $p^S$ rise, while
capital accumulation lowers $F_K$, so the inequality is eventually reversed. Single crossing is not
guaranteed in general, because productivity growth raises $F_K$ alongside $V$ --- the same
non-monotonicity recorded in Section~\ref{sec:sp:opt:results}. The date $t_R$ at
which~\eqref{eq:sp:phases:threshold} first holds with equality is an outcome of the model, not a
parameter, and is the natural definition of the onset of the circular economy in this framework.

\begin{remark}[The trigger is scale-free]\label{rem:sp:scalefree}
It is tempting to argue that recycling starts late because the marginal product of recycling capital
is low when the flow $\varpi W$ of recyclable material is small. Under the constant-returns
form~\eqref{eq:sp:setup:recycling-crs} that argument is wrong: $\mathcal R_{K^R}=a'\left(x\right)$
depends on the two inputs only through their ratio, so a small waste flow calls for a small $K^R$ at
the \emph{same} $x$, not for a low marginal product. Condition~\eqref{eq:sp:phases:threshold} contains
no quantity at all. Whether recycling occurs is decided entirely by prices and by the curvature of
$a\left(\cdot\right)$ at the origin; the scale of the economy enters only in setting how much capital
the chosen ratio requires. Earlier drafts of this model made the scale argument, and it does not
survive the rewrite of the recycling technology in constant-returns form.
\end{remark}


\subsection{Results and key mechanisms}\label{sec:sp:opt:results}

\emph{To be written once the first-order conditions and costate equations are complete.} Items to work out are recorded here as they come up.

\begin{itemize}
\item \textbf{Signs along a transition from a linear to a circular economy.} On a path where $p^W$ starts positive and turns negative once waste is chiefly a feedstock, \eqref{eq:sp:costate:M} makes $p^M$ turn negative strictly before $p^W$ does and $\zeta=\sigma\left(p^W-p^M\right)$ strictly after, so the three signs cross in a fixed order rather than together. \emph{To be stated and proved properly.} Part~(ii) of Proposition~\ref{prop:sp:phases} supplies the starting point: $p^W>0$ throughout the semi-linear phase, so any sign change occurs strictly after $t_R$.

\item \textbf{The transition as reserves deplete and productivity grows.} The starting point is an accounting identity rather than a conjecture. Rearranging~\eqref{eq:sp:setup:ledger-collected}, a materially stationary economy with no displaced mass satisfies
\begin{align}
R &= W = \frac{N}{1-a\varpi},
\label{eq:sp:multiplier-identity}
\end{align}
so \emph{recycling is a multiplier on extraction, not a substitute for it}: the circularity multiplier scales up the material throughput obtainable from a given virgin flow, but it cannot generate throughput from nothing.

Now consider depletion with technical progress. As $S$ falls, both $C^N_N$ and the rent $p^S$ rise, so~\eqref{eq:sp:sum:N} forces $N$ down; meanwhile the gross material value $\Psi$ rises on two counts, scarcity and productivity. Rising $\Psi$ raises the treatment share through~\eqref{eq:sp:sum:varpi} and pushes $p^W$ down through~\eqref{eq:sp:sum:W}, so the sign-change path of the preceding item is not a hypothetical: it is what depletion generates.

The capital margin is where we should expect the interesting non-monotonicity. Productivity growth raises $F_K$ as well as $\Psi$, and~\eqref{eq:sp:sum:KR} sets $a'\left(\Psi-p^W+d^Wp^P\right)=F_K\left(1-\zeta\Omega^Y\right)$, so the opportunity cost of putting capital into recycling rises alongside the reward. \emph{Conjecture:} the treatment share $\varpi$ rises early, since it is governed by the cost function $c^T$ alone and responds to $\Psi$ directly, while the capital ratio $K^R/K$ rises only later, once the growth of $\Psi$ outpaces that of $F_K$. Material throughput $R$ would then rise while the multiplier grows faster than $N$ falls, and decline thereafter. \emph{To be checked numerically; the ordering of the two margins is the substantive claim.} Part of the ordering is no longer a conjecture: Proposition~\ref{prop:sp:phases} establishes that $\varpi>0$ always while $K^R=0$ on an interval whenever~\eqref{eq:sp:phases:threshold} holds, so treatment strictly precedes recycling. What remains open is the behaviour of the two margins \emph{after} $t_R$, and whether~\eqref{eq:sp:phases:threshold} can be crossed more than once.

\item \textbf{Steady state, and whether growth survives depletion.} Exploration eventually ceases, since $p^X$ accumulates without bound in~\eqref{eq:sp:sum:X}, after which $S$ declines and $N\to0$. The question is whether consumption can nonetheless be sustained, and~\eqref{eq:sp:multiplier-identity} shows the answer runs through a channel absent from the classical exhaustible-resource literature: $R$ stays bounded away from zero even as $N\to0$, provided $1-a\varpi$ falls at the same rate. Since $a\to1$ requires $x=K^R/T\to\infty$ and capital is accumulable, \emph{capital substitutes for extraction inside the materials loop}, not only for materials in the production function.

Whether this suffices is a tail condition on the recycling technology. If $1-a\left(x\right)$ decays exponentially, sustaining constant $R$ against exponentially declining $N$ requires $x$ to grow only logarithmically and hence $K^R$ only linearly in time --- cheap, and consistent with sustained consumption on modest terms. If instead $1-a\left(x\right)\sim\gamma x^{-\psi}$, then $K^R$ must grow exponentially at rate $g/\psi$ where $g$ is the rate of decline of $N$, which has to be funded by productivity growth. \emph{Conjecture:} the long run is governed by the tail behaviour of $a\left(\cdot\right)$ rather than by the elasticity of substitution between capital and materials in $F$, which is where the classical condition locates it. It is worth recording that the assumption $\lim_{x\to\infty}a=1$ in~\eqref{eq:sp:setup:recycling} carries this entire result: an asymptote $\bar a<1$ caps the multiplier at $\left(1-\bar a\varpi\right)^{-1}$ and forces $R\to0$ with $N$.

The pollution side admits an exact statement. Decomposing~\eqref{eq:sp:setup:Xi} and substituting~\eqref{eq:sp:setup:ledger-collected},
\begin{align}
\Xi &= \underbrace{\left(1-d^W\right)\left(1-\varpi\right)W}_{\text{untreated waste}}+\underbrace{d^W\left(N-\dot M^K+\mathcal N\right)}_{\text{leakage on the virgin flow}},
\label{eq:sp:emissions-decomposition}
\end{align}
an identity requiring no stationarity assumption. In a materially stationary economy with full treatment it collapses to $\Xi=d^WN$: emissions are proportional to \emph{extraction}, not to output or to waste generation. Zero emissions therefore require zero extraction, hence perfect circularity, hence unbounded recycling capital. \emph{The identity is proved; its consequences for the existence and uniqueness of a steady state are not.}
\end{itemize}
